\section{Integral invariant derivatives at the counted height}\label{fg:sec:derivatives}
The height tests have a differential form on the integral Stienstra
lift. We prove the coefficientwise estimates for arbitrary lifts over
an unramified Witt ring; a specified Honda model is not a hypothesis.
These estimates occur in the earlier manuscript
\citet[\S\S2--3]{fg-invariant-derivatives}. The proofs here also give the
precise next-height test needed for the counting family.

Let $R=W(k)$ for a perfect field $k$ of characteristic $p$, and let
$\mathcal F$ be a one-dimensional commutative formal-group law over $R$.
Its normalized invariant derivation and logarithm are
\[
 D=v(T)\frac d{dT},\qquad v(T)=\partial_2\mathcal F(T,0)\in1+TR[[T]],
 \qquad L(T)=\int_0^T v(s)^{-1}\,ds.
\]
Associativity shows that $L(\mathcal F(T,U))=L(T)+L(U)$.
For the Stienstra coordinate, equation~\eqref{fg:eq:stienstra} gives
\[
 D=\left(\sum_{m\ge1}b_mT^{m-1}\right)^{-1}\frac d{dT}.
\]
For an integer $h\ge1$ put
\[
 q=p^h,\qquad c_h=\frac1{p-1}-\frac1{p^h-1},\qquad
 E_h(a)=\sum_{\substack{j\ge1\\h\nmid j}}
                    \left\lfloor\frac a{p^j}\right\rfloor.
\]

\begin{theorem}[Divisibility of invariant derivatives]\label{fg:thm:derivatives}
If the reduction of $\mathcal F$ has height at least $h$, then for every
$a\ge0$ and $f\in R[[T]]$,
\begin{equation}\label{fg:eq:derivative-bound}
 D^af\in
 p^{\max\{0,E_h(a),\lceil c_ha-(h-1)\rceil\}}R[[T]].
\end{equation}
If its height is exactly $h$, then for $j=hr+s$, $0\le s<h$,
\[
 v_p\bigl(D^{p^j}(T^{p^s})(0)\bigr)=E_h(p^j),\qquad
 v_p\bigl(D^{q^r}T(0)\bigr)=c_h(q^r-1).
\]
Thus $c_h$ is the optimal uniform asymptotic slope at exact height $h$.
The error $h-1$ in the affine bound is optimal for every fixed normalized
group satisfying the height lower bound.

At infinite height, $L$ and $L^{-1}$ are integral and
\[
 D^aR[[T]]\subseteq p^{v_p(a!)}R[[T]].
\]
This exponent is attained. The slope $1/(p-1)$ has no bounded uniform
affine error in this infinite-height case.
\end{theorem}

\begin{proof}
We first prove the affine estimate. By Lemma~\ref{fg:lem:height-log},
the coefficients of $L$ below degree $q$ are integral. The new coordinate
\[
 z=T+\sum_{2\le m<q}[T^m]L(T)\,T^m
\]
has an integral compositional inverse. In this coordinate the logarithm
is $z+bz^q+O(z^{q+1})$, with $pb\in R$. A strict integral coordinate
change conjugates $D$ and preserves every ideal $p^eR[[T]]$, so it
suffices to work in this coordinate.

Define integral translation operators $\Delta_i$ by
\[
 f(\mathcal F(z,u))=\sum_{i\ge0}(\Delta_i f)(z)u^i.
\]
Logarithmic addition gives
$\sum_i\Delta_i u^i=\exp(L(u)D)$. Coefficient comparison yields
\[
 \Delta_i=D^i/i!\quad(i<q),\qquad
 \Delta_q=D^q/q!+bD.
\]
Set $B=v_p(q!)-1=c_h(q-1)$ and write $q!=p^{B+1}w$, with $w$ a unit.
Then
\begin{equation}\label{fg:eq:operator-recursion}
 D^q=p^BuD+p^{B+1}A,\qquad u=-wpb\in R,
 \qquad A=w\Delta_q,
\end{equation}
where $A$ preserves the integral function ring.

For $0\le a<q$, the equality $D^a=a!\Delta_a$ gives factorial
divisibility. Writing $a=\sum_{i=0}^{h-1}a_i p^i$, Legendre's formula
gives
\[
 v_p(a!)-c_ha
 =\sum_{i=0}^{h-1}a_i
       \left(\frac{p^i}{q-1}-\frac1{p-1}\right)\ge1-h.
\]
Every parenthesized coefficient is nonpositive, so the minimum occurs
when all digits are $p-1$, at $a=q-1$.
For $a\ge q$, apply \eqref{fg:eq:operator-recursion} to $D^{a-q}f$.
Induction bounds the first term by $c_ha-(h-1)$ and the second by
that quantity plus $1-c_h\ge0$. This uses only that $A$ preserves
integral ideals; it needs no commutation of $A$ with $D$.
Integer-valued valuations and integrality prove the affine estimate.
For every normalized coordinate,
\[
 D^{q-1}(T^{q-1})(0)=(q-1)!,\qquad
 v_p((q-1)!)=c_h(q-1)-(h-1).
\]
This proves optimality of the error independently of the slope argument.

For the floor-sum estimate, adjoin $\pi$ with $\pi^{q-1}=p$ and work
in the discrete valuation ring $R'=R[\pi]$. The rescaled law
\[
 \mathcal F_\pi(X,Y)=\pi^{-1}\mathcal F(\pi X,\pi Y)
\]
is integral, and every coefficient of its $p$-series is divisible by
$p$. Below degree $q$ this follows from the height hypothesis; in
degree $m\ge q$ it follows from the factor $\pi^{m-1}$.
The recursion \eqref{fg:eq:log-recursion}, in every degree, therefore
makes its logarithm integral over $R'$, and its inverse is integral.
If $E=L^{-1}=\sum e_mT^m$, uniqueness gives
\[
 E_\pi(T)=\pi^{-1}E(\pi T),\qquad
 v_p(e_m)\ge-\frac{m-1}{q-1}.
\]
This assertion uses divisibility by $p$ of the entire rescaled
$p$-series, rather than only reduction modulo the ramified uniformizer.

Logarithmic differentiation gives
\[
 D^af(0)=a![T^a]f(E(T)).
\]
A product of $b\ge1$ exponential coefficients with total degree $a$
has valuation at least $-(a-b)/(q-1)$. Since the resulting coefficient
lies in the unramified field $R[1/p]$, its valuation is an integer.
Consequently
\[
 v_p(D^af(0))\ge
 v_p(a!)-\left\lfloor\frac{a-1}{q-1}\right\rfloor
 \quad(a\ge1).
\]
The same estimate holds coefficientwise: apply the evaluation argument
in $u$ to each coefficient in $T$ of $f(\mathcal F(T,u))$, using
$D_u^a f(\mathcal F(T,u))|_{u=0}=D_T^af(T)$.

Now write $a=\sum_j a_jp^j$. Compose these coefficientwise estimates
for the powers $(D^{p^j})^{a_j}$. The finite identities
\[
 \sum_j a_jv_p((p^j)!)=v_p(a!),\qquad
 \sum_j a_j\left\lfloor\frac{p^j-1}{q-1}\right\rfloor
       =\sum_{r\ge1}\left\lfloor\frac a{q^r}\right\rfloor
\]
give $v_p(D^af)\ge E_h(a)$. Together with the affine estimate this
proves \eqref{fg:eq:derivative-bound}.

Suppose the height is exactly $h$, and let $\gamma_q$ be the coefficient
of $T^q$ in $[p]_{\mathcal F}$. It is a unit. Modulo $\pi$ the rescaled
$p$-series divided by $p$ is $T+\overline\gamma_qT^q$.
Indeed, degrees below $q$ other than one have an additional positive
power of $\pi$, and degrees above $q$ do also. Divide the logarithm
recursion by $p$. All its terms involving the $i$-th power of the
$p$-series for $i\ge2$ vanish modulo $\pi$. Hence
\[
 \overline L_\pi(T)=T+\overline\gamma_qT^q,\qquad
 \overline E_\pi(T)=\sum_{r\ge0}\beta_rT^{q^r},
\]
where $\beta_0=1$ and $\beta_r=-\overline\gamma_q\beta_{r-1}^q$.
Every $\beta_r$ is nonzero. For $j=hr+s$, the coefficient of $T^{p^j}$
in $\overline E_\pi(T)^{p^s}$ is therefore nonzero. Undoing the scaling
gives
\[
 v_p\bigl([T^{p^j}]E(T)^{p^s}\bigr)
 =-\frac{p^j-p^s}{q-1}
 =-\left\lfloor\frac{p^j-1}{q-1}\right\rfloor.
\]
Multiplication by $(p^j)!$ proves both witness equalities. Letting
$r$ tend to infinity in the witness with $f=T$ excludes every larger
slope with bounded, or even sublinear, error for this fixed group.

At infinite height every coefficient of $[p]$ belongs to $pR$.
The logarithm recursion proves $L\in R[[T]]$ in every degree, and
its inverse is integral. In its coordinate, $D=d/dL$, whose $a$-th
iterate is coefficientwise divisible by $a!$. The function $L(T)^a$
attains $a!$. Finally
\[
 v_p(a!)=\frac a{p-1}-\frac{s_p(a)}{p-1}.
\]
At $a=p^r-1$ the digit error is $r$, which is unbounded. This proves
the infinite-height assertions.
\end{proof}

\begin{corollary}[A differential test for the next height]\label{fg:cor:derivative-height}
For a Stienstra lift whose reduction has height at least $r$, put $q=p^r$.
Then
\[
 \frac{D^qT(0)}{q!/p}\equiv
       -\frac{b_q}{p^{r-1}}\pmod p,
\]
and both quotients are integral. Moreover
\[
 h(X)\ge r+1
 \iff D^qT(0)\in q!R
 \iff D^qR[[T]]\subseteq q!R[[T]].
\]
For $r=2$, the valuation of $D^{p^2}T(0)$ is exactly $p$ at height
two and at least $p+1$ at height at least three.
\end{corollary}

\begin{proof}
Write $L=T+\sum\ell_mT^m$ and $E=T+\sum e_mT^m$.
Lemma~\ref{fg:lem:height-log} and formal inversion give
$e_m\in R$ for $m<q$, and $e_q+\ell_q\in R$. Since
$D^qT(0)=q!e_q$, this gives the signed congruence and the pointwise
equivalence. At height at least $r+1$, the floor sum in
Theorem~\ref{fg:thm:derivatives}, with parameter $r+1$ and order $q$,
is $v_p(q!)$. This proves the whole-operator inclusion. The converse
follows by testing $T$ at zero. Finally $v_p((p^2)!)=p+1$.
\end{proof}

For the counting family, $S=0\bmod p$ gives height at least two.
Every integral formal function on its Stienstra lift consequently
satisfies
\[
 D^af\in p^{\max\left\{0,
    \sum_{j\ge0}\left\lfloor a/p^{2j+1}\right\rfloor,
    \left\lceil pa/(p^2-1)-1\right\rceil\right\}}R[[T]].
\]
In particular $\overline D^{\,p}=0$. This reduced operator vanishes
at every height at least two, including infinity; the next distinction
is the additional factorial divisibility in
Corollary~\ref{fg:cor:derivative-height}.

\subsection{The base and the Honda-dual hypothesis}
The unramified base and the scope of the height hypothesis are essential.
Over $\Z_p[[u]]$ the law $X+Y+uXY$ has
\[
 D=(1+uT)\frac d{dT},\qquad D^aT=u^{a-1}(1+uT)\quad(a\ge1).
\]
At $u=0$ its reduction is additive, while $D^pT$ is not divisible by
$p$ over the entire parameter ring. If $u=\pi$ in a ramified discrete
valuation ring with $v_\pi(p)=e>p-1$, the residue group is again additive,
but $v_p(D^pT(0))=(p-1)/e<1$. These computations explain exactly why
height at one closed point, or reduction only modulo a ramified
uniformizer, does not give the theorem's coefficientwise $p$-ideals.

The directional result in \citet[\S6]{fg-invariant-derivatives} concerns
the Cartier dual of the specified Honda lift with logarithm
$\sum_{r\ge0}T^{p^{hr}}/p^r$. That model has $h-1$ specified primitive
directions, with affine slopes $p^i/(p^h-1)$ and error one for
$1\le i<h$. An integral identification of that model and its directions
is an additional hypothesis. Neither the counting construction nor
Stienstra's formula provides it, so that directional result is not used
in the preceding proofs.

Equal height need not even give an isomorphism over the prime field.
For the Fermat cubic over $\Z_5$,
\[
 b_5=0,\qquad b_{25}=24!/(8!)^3=9465511770,\qquad
 b_{25}/5=4\pmod5.
\]
Its height-two leading $p$-series coefficient is four. The standard
Honda law just specified has coefficient one. An isomorphism over
$\F_5$ would multiply that coefficient by $u^{1-25}=1$ for
$u\in\F_5^\times$, which cannot identify them. In particular their
integral lifts over $\Z_5$ are not isomorphic. This is an obstruction
to inferring the model identification, rather than an assertion that
every numerical directional estimate fails in every other model.
